\section{Introduction}\label{sec:introduction}

Critical FK--Ising interfaces converge to chordal
$\mathrm{SLE}_{16/3}$. Moving the temperature away from criticality
on the scale of the mesh retains a macroscopic mass and leads to a
different interface law. The term \emph{massive SLE} describes this
expected deformation, but does not by itself specify a continuum
process or establish that lattice approximations have a unique
limit. The aim of this paper is to provide both the limit and a
well-posed continuum characterization for the thermal FK--Ising
model.

The off-critical interface program was formulated by Makarov and
Smirnov \cite[Problems~A--B and Section~2.4]{MS}: find the scaling
limit of the relevant observable, derive a driving equation, and
prove uniqueness and convergence of the interfaces. For FK--Ising,
the fermionic boundary value problem makes the second step
substantially different from the massive harmonic examples in that
work. At criticality, the observable approach originates in
Smirnov's work \cite{Smirnov}; isoradial observables were developed
by Chelkak and Smirnov \cite{CS}, and convergence of the Ising and
FK--Ising interfaces was proved in \cite{CDCHKS}. Park
\cite{Park} established the massive two-point fermionic observable
and uniform near-critical rectangle crossings on isoradial
lattices. These supply the observable and crossing inputs; identifying a unique
interface law still requires the stochastic analysis carried out here. The general crossing-to-Loewner compactness
theory of Kemppainen and Smirnov \cite{KS} supplies the passage
from uniform conditional crossing bounds to subsequential curves
and drivers.

We prove that the thermal interfaces converge along the full
sequence for every bounded simply connected marked domain, both
signs of the mass, varying uniformly angle-bounded isoradial
lattices, and every nome sequence with the prescribed asymptotic
ratio. The approximating domains need only converge in the marked
Carath\'eodory sense; their diameters may diverge. Curves are
compared after normalized transfer to the closed disc, so no
regularity of physical boundary impressions is required. This
settles the scaling-limit and uniqueness problem for the thermal
FK--Ising interfaces in the formulation of
Theorem~\ref{thm:main}.

\Needspace{16\baselineskip}
The limiting equation is particularly concrete. For $m>0$, in the
active physical domain let $h$ be the bounded solution of
\[
 \Delta h=-4m|\grad h|
\]
with value $1$ on the primal-wired bank and $0$ on the dual-wired
bank. These map respectively to the negative and positive real axes
in the current half-plane chart centered at the tip. Write $\psi$ for the
inverse conformal map and $\mu=m|\psi'|$. For positive mass the
driving equation is
\[
 \dd W_t=\sqrt{16/3}\,\dd B_t+\frac{2\pi}{3}C_t\dd t,
 \qquad
 C_t=\frac4\pi\int_{\Hh}\Ima\frac{-1}{v}\,
       \mu_t(v)|\grad(h_t\circ\psi_t)(v)|\,\dd A(v).
\]
The variance, the absolute continuity of the drift, and its finite
energy are conclusions of the proof. They are not assumed in the
compactness argument. Weak uniqueness holds in the finite-energy
class; it yields a definition independent of a lattice subsequence.
The negative-mass law follows by duality and reversal.

There are three main technical points. First, we obtain a
moving-mark energy estimate before knowing that the driver is a
semimartingale. Conditional crossing bounds give restarted control
of the angle seen from a bulk point. Reflection near an avoided
boundary interval and a small-power contraction in the point's
Loewner clock then yield
\[
 \E\int_0^\infty
   \left(\int_{\Hh}\frac{m|\psi_t'(v)|}{|v|^2}\,\dd A(v)\right)^2
   \dd t<\infty.
\]
The proof handles small initial angles and infinite clock lifetime
at a finite swallowing time in capacity parametrization. Its dependence on crossing
geometry rather than on an already known stochastic differential
equation makes the estimate useful for identifying new interface
limits.

Second, a height cutoff justifies stochastic Green-potential
calculus up to swallowing. A common clock gives simultaneous densities for capacity,
the variation of the driver's finite-variation part, and spatial
martingale brackets. Averaged near-tip expansions then show that
the quadratic variation of the driver's martingale part and its
finite-variation measure are absolutely continuous with respect
to capacity time, and identify their densities in the equation above. This avoids assuming a boundary
value for the martingale bracket field at the moving tip. A
stopped-Girsanov argument subsequently proves uniqueness for the
measurable drift functional, including a precise treatment of its
energy stopping times.

Third, kernel convergence alone does not bound the physical size
of the lattice domains. We remove remote pieces by a boundary
cluster exploration that leaves an ordinary bounded Dobrushin
problem. Restricted length--area estimates control the normalized
images of remote crosscuts. A hierarchy of mandatory gates gives
a summable failure bound even when there are arbitrarily many
entrances to distant regions. This localization preserves the full
disc curve topology and avoids using near-critical rectangle
bounds at diverging physical scales.

The proof is organized by these dependencies. Section~\ref{sec:model}
states the model and continuum prescription.
Section~\ref{sec:geometry} constructs the subsequential martingale
fields. Sections~\ref{eng:section} and \ref{cal:section} establish the
energy estimate and identify the driver. Section~\ref{sec:uniqueness} proves uniqueness and the
bounded-approximation result, and
Section~\ref{sec:localization} completes the arbitrary-domain
and negative-mass cases.
